V matematické logice rozlišujeme tři zásadní pojmy: \textbf{jazyk},
\textbf{teorii} a \textbf{model}.

\begin{definition}[Jazyk]
  \label{def:jazyk}
  Jazykem myslíme soubor symbolů (abecedu) spolu s gramatikou -- pravidly pro
  spojování symbolů (tj. určení, které posloupnosti symbolů jsou správné a které
  ne).
\end{definition}

\begin{parable}
  \label{par:jazyk}
  O \emph{jazyce} ve smyslu matematické logiky lze přemýšlet jako o češtině --
  jenom některé posloupnosti písmen jsou českými slovy (např. \emph{gyufj} není,
  zatímco \emph{kočka} je) a pouze některé posloupnosti slov jsou gramaticky
  správné věty.

  \textbf{Pozor!} \emph{Gramaticky správné} věty ještě nemusejí
  \uv{dávat smysl}. Gramatika pouze určuje, jak lze stavět písmena do slov a
  slova do vět, nikoli které z těchto staveb mají \uv{věcný význam}. Čili, třeba
  věta: \uv{Uzený pes rozkládal v hořčičné periferii,} je gramaticky správně,
  ačkoli spíše nedává smysl.
\end{parable}

Jazyku přiřazuje význam teprve tzv. \emph{teorie}.

\begin{definition}[Teorie]
  \label{def:teorie}
  Ať $L$ je nějaký \hyperref[def:jazyk]{jazyk}. Teorií v jazyce $L$ rozumíme
  soubor gramaticky správných posloupností symbolů v jazyce $L$.

  Intuitivně vzato je teorie výčet vět v jazyce $L$, které \uv{dávají smysl}.
\end{definition}

\begin{parable}
  \label{par:teorie}
  Je-li jazyk $L$ čeština, pak jednou jeho možnou teorií je soubor všech vět,
  které nám dávají smysl, i když nemusejí být nutně fakticky správné. Například
  již vyřčená věta \uv{Uzený pes rozkládal v hořčičné periferii,} by součástí
  takové teorie nebyla, ale třeba \uv{Mám rád modrou barvu,} již ano.
\end{parable}

Konečně, \emph{model} je struktura, která vznikne, když všechny věty z teorie
označíme za \emph{pravdivé} nebo \emph{fakticky správné}.

\begin{definition}[Model]
  \label{def:model}
  Ať $L$ je jazyk a $T$ teorie v jazyce $L$. \emph{Modelem} teorie $T$ je
  struktura, v níž jsou všechny věty (či formálněji všechny gramaticky správné
  posloupnosti symbolů) teorie $T$ pravdivé.
\end{definition}

\begin{parable}
  Označíme-li každou gramaticky správnou českou větu za pravdivou, tak dostaneme
  poněkud nesmyslný model (s každou větou v něm bude platit i její opak).
  Představme si ideální situaci, kdy jsme všeznalí a přesně víme, které věty
  jsou pravdivé a které lživé. Když všechny \emph{pravdivé} (a žádné
  \emph{lživé}) věty zahrneme do teorie $T$ jazyka českého, pak jeho modelem
  bude přesně náš svět -- tedy struktura, ve které platí každá věta, o níž při
  své všeznalosti víme, že je pravdivá.
\end{parable}

Nyní uvedeme \emph{jazyk} výrokové a predikátové logiky. Spolu s pár dalšími
symboly tvoří tyto jazyk \emph{teorie} množin, teorie, jejímž jedním modelem je
většina moderní matematiky.

\begin{figure}[ht]
  \centering
  \begin{tikzpicture}
    \begin{scope}
      % Ellipse with label
      \node[ellipse,draw=FireBrick,minimum width=4cm,minimum
      height=2cm,thick] (ej) at (0,0) {};
      \node at (90:1.4) {\clr{\textsf{Jazyk}}};

      % Alphabet
      \node at (0,0.5) {a};
      \node at (0.6, 0.4) {z};
      \node at (1, -0.3) {j};
      \node at (-0.7, -0.2) {.};
      \node at (-0.2, 0.1) {?};
      \node at (0.4, -0.7) {f};
      \node at (-0.7, 0.4) {Y};
      \node at (0.2, 0) {/};
      \node at (-0.1, -0.5) {W};
    \end{scope}
    \draw[thick,-Latex] (2.1,0) -- (2.9,0);
    \begin{scope}[xshift=5cm]
      % Ellipse with label
      \node[ellipse,draw=RoyalBlue,minimum width=4cm,minimum
      height=2cm,thick] (et) at (0,0) {};
      \node at (90:1.4) {\clb{\textsf{Teorie}}};

      % Sentences
      \node at (0, 0.3) {Slunce zapadá.};
      \node at (0, -0.3) {Jelen se pase.};
    \end{scope}
    \draw[thick,-Latex] (7.1,0) -- (7.9,0);
    \begin{scope}[xshift=10cm]
      % Ellipse with label
      \node[ellipse,draw=ForestGreen,minimum width=4cm,minimum
      height=2cm,thick] (em) at (0,0) {};
      \node at (90:1.4) {\clg{\textsf{Model}}};

      % Picture
      % TODO get picture
    \end{scope}
  \end{tikzpicture}
  \caption{Vztah mezi jazykem, teorií a modelem.}
  \label{fig:jazyk-teorie-model}
\end{figure}
